% !PRINTED 524
\weberpage{524}

By \S.\,142, 4., if $\omega$ is a given integral functional, $E$ a unit, and $\varphi$ an integral function in $\Omega$, we can set:

\weq{1}{E\omega = \varphi(x, y, \ldots).}

If, on the other hand, we decompose $\omega$ into its prime factors, then by 1. we may assume these prime factors to be free of the variables $x, y, \ldots$, and if we again form the product of these prime factors, we obtain a number $\omega_1$ associated with $\omega$, which is free of the variables $x, y, \mathinner{.\,.}$.

If we order the quotient $\varphi : \omega_1$, which is an integer (indeed a unit), according to the variables $x, y, \ldots$, then we conclude from \S.\,142, 3. that all the coefficients of the function $\varphi$ are divisible by $\omega_1$, and consequently also by $\omega$ and $\varphi$.

If, on the other hand, all the coefficients of $\varphi$ are divisible by some factor $\delta$ in $\Omega$, then $\varphi$ too is divisible by $\delta$, and from this it follows:

2. \gesperrt{An integral function $\varphi$ with integers as coefficients is the greatest common divisor of all its coefficients.}

The integer $\varphi$ therefore passes, by \S.\,139, into an associated number when the individual powers and products of the variables are each replaced by a single variable. From this it also follows as a corollary that every functional $\omega$ passes into an associated one when the variables are in any way denoted differently.

We now have to prove the following important theorem:

3. \gesperrt{If $\omega$ is an arbitrary integral functional, one can choose a number $\alpha$ divisible by $\omega$ in such a way that the quotient $\alpha : \omega$ is relatively prime to an arbitrarily given functional $\mu$.}

We first prove that there is an integer $\alpha$ which is divisible by $\omega$, but not by $\omega\pi$, where $\pi$ is an arbitrary prime functional.

For if, following \S.\,142, 4., we form an integral function $\varphi$ associated with $\omega$, then the coefficients of this function are indeed all divisible by $\omega$, but not all divisible by $\omega\pi$, since otherwise $\varphi$, and hence $\omega$ itself, would be divisible by $\omega\pi$, which is not possible. Among the coefficients of $\varphi$ there is therefore at least one which has the required property.
